\documentclass[10pt,twocolumn]{article}
\usepackage[scaled]{helvet}
\renewcommand\familydefault{\sfdefault}
\usepackage[T1]{fontenc}
\usepackage[margin=0.7in]{geometry}
\usepackage{titlesec}
\usepackage{enumitem}
\usepackage{parskip}
\setlength{\columnsep}{0.24in}
\setlist[itemize]{leftmargin=1.2em, topsep=0.2em, itemsep=0.18em}
\titleformat{\section}{\large\bfseries}{\thesection}{1em}{}

\begin{document}
\title{\vspace{-1.6cm}AWS VPC (Virtual Private Cloud)}
\author{AWS ML Study Notes}
\date{}
\maketitle
\small

\section*{Overview}
A VPC is your logically isolated network in AWS. It defines the IP range, subnets, routing, and traffic boundaries for the compute and managed services you place inside it.

ML systems frequently handle sensitive data and expensive infrastructure, so network design matters. A good VPC layout reduces exposure, controls egress, and separates public entry points from private processing layers.

\section*{Core Concepts}
\begin{itemize}
\item Subnets divide the VPC into smaller network zones, usually across multiple Availability Zones for resilience. Public subnets route to an internet gateway, while private subnets usually route outbound through a NAT or VPC endpoints.
\item Security groups are stateful virtual firewalls attached to resources, while network ACLs are stateless controls applied at the subnet boundary.
\item Route tables, VPC endpoints, peering, Transit Gateway, and DNS settings determine how traffic moves between services, accounts, and external networks.
\end{itemize}

\section*{How It Works}
\begin{itemize}
\item You choose a CIDR range, create subnets, attach route tables, define security controls, and then launch resources into the appropriate subnets.
\item Private workloads such as databases, internal APIs, or SageMaker jobs can access AWS services through VPC endpoints without traversing the public internet.
\item Ingress is usually limited to controlled entry points such as an Application Load Balancer, API Gateway integration, or VPN connection from trusted environments.
\end{itemize}

\section*{AI and ML Relevance}
\begin{itemize}
\item A common ML pattern is private training or inference infrastructure that reads data from S3 through endpoints and writes logs and metrics without needing public IP addresses.
\item Network isolation is also important for compliance and for protecting proprietary models and feature data from accidental public exposure.
\end{itemize}

\section*{Operational Notes}
\begin{itemize}
\item Plan CIDR blocks early. Renumbering a VPC later is painful, especially when multiple accounts or peered networks are involved.
\item Use multiple Availability Zones for resilience, and capture VPC Flow Logs when network troubleshooting or auditability matters.
\item Do not overexpose resources by putting everything in public subnets. Public accessibility should be the exception, not the default.
\end{itemize}

\section*{What To Remember}
\begin{itemize}
\item Security groups normally express allowed flows better than broad subnet ACL rules.
\item Public subnet does not mean the instance must be public. Internet reachability still depends on routing and public addressing.
\item For secure ML platforms, private subnets plus endpoints are usually the baseline design.
\end{itemize}

\end{document}
